% Paste inside the existing \matmethods{...} block of the PNAS template.
\subsection*{Ingredient network and permutation analysis}
We used 26,727 Food.com recipes assigned to exactly one of nine cuisine tags in the project dataset, derived from the Food.com Recipes and Interactions collection \cite{Majumder2019}. Ingredient names were lowercased and mapped with 72 manually reviewed aliases. We excluded 20 generic names (including salt, pepper, water, and unspecified oil), ingredients present in fewer than three recipes, and ingredients present in more than 35\% of recipes. The remaining binary recipe-by-ingredient matrix contained 3,877 ingredient types. For recipes $i$ and $j$, similarity was the Jaccard index $J(i,j)=|I_i\cap I_j|/|I_i\cup I_j|$ over retained ingredient sets. We required at least two shared ingredients and $J(i,j)\geq0.18$, then retained the union of each recipe's 15 highest-similarity neighbors. Edge weights equaled Jaccard similarity; isolated recipes remained in the graph.

The global within-cuisine statistic was the fraction of total edge weight joining equal cuisine labels. Its reference distribution came from 199 permutations of labels over the fixed graph, preserving edge weights and cuisine counts. We detected communities using weighted Louvain optimization and compared nonisolated community assignments with cuisine labels using adjusted mutual information; 199 label permutations provided its reference distribution. For each distinct cuisine pair $(a,b)$, we divided its total edge weight by the total weight of all cross-cuisine edges. We compared this conditional pair share with the mean from 999 fixed-graph label permutations. These pair comparisons describe the network and are not formal tests of individual pairs. We also varied the minimum Jaccard threshold (0.12, 0.18, and 0.24) and manually reviewed 12 of 60 candidate edges from three selected pairs. The manual review was a qualitative spot check, not an estimate of labeling-error prevalence.

As a stricter ingredient-level reference, we made 49 independent randomizations of the binary recipe--ingredient matrix within each cuisine using Curveball trades, with ten attempted trades per recipe. A trade exchanges ingredients unique to two recipes while retaining their shared ingredients. It preserves exactly the number of retained ingredients in every recipe and the number of recipes containing every retained ingredient within each cuisine. For each randomized matrix, we rebuilt the Jaccard/top-15 network and recalculated the within-cuisine weight share and three illustrative cross-cuisine pair shares. The observed graph was reproduced edge for edge before randomization; row and cuisine-specific column sums were verified after every replicate. Empirical upper-tail $P$ values use a one-count correction and are exploratory for pairs selected after inspecting the observed network. Five additional randomizations with twice as many attempted trades per recipe checked the stability of the reference medians.

We recovered non-cuisine course tags from the original Food.com \texttt{RAW\_recipes.csv}, matching all 26,727 recipe IDs to unique raw rows. We assigned each recipe one primary dish type by applying the following precedence to its tags: beverage, dessert, soup/stew, salad, bread, breakfast/brunch, appetizer/snack, side dish, condiment/sauce, and main dish; recipes without these tags were assigned other/unclear. The assignment yielded 11 types, including 1,628 other/unclear recipes; 5,679 recipes had multiple course tags, so the primary type is a simplification. We made another 49 independent Curveball randomizations, now trading ingredients only among recipes in the same cuisine and primary-type stratum (99 nonempty strata). These trades preserved each recipe's ingredient count and each ingredient's frequency within every stratum. We verified all margins, rebuilt the network, and compared the same statistics with the observed graph and the earlier cuisine-only randomizations. A title-word classification agreed with the tag-assigned type for 3,662 of 4,201 explicitly named recipes; this check is not an accuracy estimate because titles are also ambiguous.

For diagnostics, we ran five additional cuisine-by-primary-type randomizations with 20 attempted trades per recipe and five randomizations grouping by the complete combination of course tags within cuisine. The latter had 823 nonempty strata, including 274 singletons. These checks compare reference medians but are too small for meaningful tail-probability estimates.

For sensitivity checks, we applied the same post-construction edge filters to the observed graph and every randomized graph: removal of edges incident to recipes explicitly named as sweets or drinks by conservative title rules; removal of edges incident to original Louvain community 76; and removal of edges with Jaccard similarity at least 0.75. Community membership was fixed from the observed graph, so its exclusion is descriptive rather than an independent test. Remaining nodes were not reconnected after filtering.
